\section{Prediction and exact spectral repair}\label{sec:repair}

The Fourier construction in Section~\ref{sec:prime-kernel} expresses
a large transform as a product of \(m\) commuting operators. Applying
them one at a time would cost \(m\) stages. In a common eigenbasis their
product acts by a single scalar at each frequency, but changing to that
basis on every input row would itself be expensive. We use a low-degree
predictor on a batch of independent rows. At each frequency it is exact
on almost all rows, and a small number of linear checks recovers the
product on the whole batch. The exceptional rows may change from one
frequency to the next.

Let \(\ell\) be prime, let \(\zeta\) be a primitive complex \(\ell\)-th
root of unity, and set \(R_\ell=\mathbb Z[1/\ell,\zeta]\). A row is an
array indexed by a finite set, with coefficients in a commutative
\(R_\ell\)-algebra. Suppose that one invertible change of coordinates
\(\mathcal S\) diagonalizes \(Q_1,\ldots,Q_m\) on this array. At frequency
\(\nu\), operator \(Q_j\) has multiplier \(\zeta^{k_j(\nu)}\), and the
exponents \(k_j(\nu)\) are given. Write \(F=\prod_{j=1}^m Q_j\).

We count two operations on a batch of \(W\) rows. A \emph{pencil stage}
combines the query powers with scalar matrices acting on the row index,
\begin{equation}\label{eq:repair-pencil}
  A_0\otimes I+\sum_{j=1}^m\sum_{a=1}^{\ell-1}A_{ja}\otimes Q_j^a,
       \qquad A_0,A_{ja}\in\operatorname{Mat}_{W\times W}.
\end{equation}
A \emph{spectral table on \(h\) rows} takes \(h\) forward conversions by
\(\mathcal S\), applies one prescribed \(h\times h\) scalar matrix at
each frequency, and takes \(h\) inverse conversions. Pointwise linear
maps combine rows at the same ordinary index. A pencil stage can
contain all \(m\) queries. The next section bounds the cost of executing
such a stage in the Fourier construction.

\needspace{4\baselineskip}
\begin{theorem}[Simultaneous spectral repair]\label{thm:repair}
For every \(m\ge2\) there are a positive batch width \(W\) and a finite
free cyclotomic extension of \(R_\ell\), depending only on \(\ell,m\),
over which a program computes \(F\) on \(W\) independent rows. It uses
\(O_\ell(\sqrt{m\log(2m)})\) pencil stages, pointwise linear maps and
one spectral table on \(2s\) rows, with \(s\le W/(2m^3)\). The table
thus uses at most \(4s\) full row conversions. Its matrix at frequency
\(\nu\) depends only on \((k_j(\nu))_{j=1}^m\). All row matrices and the
finite list of spectral matrices are effectively constructible for fixed
\(\ell,m\). The identities specialize to any commutative \(R_\ell\)-algebra
after base change along the finite free extension.
\end{theorem}

\subsection{Predicting the product from an integer count}

Assign each row a shift \(\beta\in(\mathbb Z/\ell\mathbb Z)^m\), using
every shift equally often. At a fixed \(\nu\), write \(q_j=\zeta^{k_j(\nu)}\)
and \(z=\prod_jq_j\) for the desired multiplier. Let \(b_j\) be the least
nonnegative residue of \(k_j(\nu)+\beta_j\), and put \(c=\sum_j b_j\)
as an ordinary integer. On a row with shift \(\beta\), it suffices to
predict \(\zeta^{c-\sum_j\beta_j}=z\). Varying the shift makes the digits
\(b_j\) uniform and independent at every fixed frequency.

For three binary queries, for example, \(c\) is the Hamming weight of
\(k+\beta\). On the window \(1\le c\le2\), the linear predictor
\((-1)^{|\beta|}(2c-3)\) gives the target sign \((-1)^{|k|}\).
The two shifts outside the window are \(000,111\) at \(k=000\), but
\(001,110\) at \(k=001\).

In general, let \(r\) be the least integer with \(2^r\ge2m\). Set
\[
 \mu=m(\ell-1)/2,\qquad B=(\ell-1)\lceil\sqrt{2mr}\rceil,\qquad
 u=\lfloor\mu-B\rfloor,\quad v=\lceil\mu+B\rceil .
\]
These endpoints are computable by integer arithmetic. Hoeffding's
inequality \cite[Theorem~2]{hoeffding} bounds the fraction of rows
outside this window at each frequency,
\[
 \Pr_\beta[c\notin[u,v]]
    \le 2\exp\!\left(-\frac{2B^2}{m(\ell-1)^2}\right)
    \le \frac{1}{2m^3}.
\]
The predictor will have degree at most the width \(D=v-u\). We use
integral interpolation so that its coefficients also make sense in
positive characteristic.

\begin{lemma}[Integral interpolation on the count]\label{lem:count-predict}
For each shift \(\beta\) there is a polynomial \(P_\beta\) over
\(R_\ell\), of degree at most \(D=v-u\) in the variables representing
the query powers \(Q_j^a\), \(1\le a<\ell\). At each frequency with
\(c\in[u,v]\) this polynomial equals the multiplier of \(F\).
\end{lemma}
\begin{proof}
The indicator for the event \(b_j=b\) is affine-linear in the query powers,
\[
 I_{j,b}=\frac1\ell\sum_{a=0}^{\ell-1}\zeta^{a(\beta_j-b)}q_j^a.
\]
For \(h>0\), the coefficient of \(T^h\) in \((1+T)^{b_j}\) is thus
represented by \(\sum_{b=0}^{\ell-1}\binom bh I_{j,b}\). Use \(1\)
for the constant coefficient. In the formal expansion
\[
 L_h=[T^h](1+T)^{-u}\prod_{j=1}^m(1+T)^{b_j}
\]
each term uses a positive power of \(T\) from at most \(h\) query
factors. It follows that \(L_h\) has a representation of degree at
most \(h\) in the query powers. The fixed series \((1+T)^{-u}\) has
integer coefficients.

Now take
\[
 P_\beta=\zeta^{\,u-\sum_j\beta_j}
                \sum_{h=0}^{D}(\zeta-1)^h L_h .
\]
On the window, \(0\le c-u\le D\). The truncated sum therefore evaluates
the whole polynomial \((1+T)^{c-u}\) at \(T=\zeta-1\) and returns
\(\zeta^{c-\sum_j\beta_j}=z\) after scaling.
\end{proof}

\subsection{Evaluating a predictor on independent rows}

A polynomial of degree \(D\) can have many monomials. To evaluate it in
\(O_\ell(D)\) pencil stages, increase the batch width and express the
polynomial times the identity as a power of one matrix of affine-linear
forms. The identity factor ensures that every coordinate of the matrix
holds an independent input.

Childs studied such linearizations of forms through generalized Clifford
algebras \cite{childs}. For homogeneous forms over a field with a
primitive root of the required order, see Bl\"aser, Eisenbud and
Schreyer \cite[Theorem~2.1]{ulrichcomplexity}. We give the construction
over a localized cyclotomic ring to obtain a pencil that specializes in
every characteristic allowed here.

\begin{lemma}[Matrix evaluation of a polynomial]\label{lem:matrix-root}
Let \(p\) be a polynomial of degree at most \(b\) over \(R_\ell\), where
\(b=\ell^a\ge2\). If \(\xi\) is a primitive \(b\)-th root of unity,
there is a square matrix \(L\) of affine-linear forms over
\(R_\ell[\xi]\) such that \(L^b=pI\).
\end{lemma}
\begin{proof}
Write \(p=\sum_{i=1}^h\prod_{r=0}^{b-1}\lambda_{ir}\), padding each
monomial with constant factors so that every \(\lambda_{ir}\) is
affine-linear. On the \(i\)-th copy of a free module with basis
\(e_0,\ldots,e_{b-1}\), put
\[
 S_i e_r=\lambda_{ir} e_{r+1\bmod b},\qquad Z e_r=\xi^r e_r .
\]
Then \(S_i^b=(\prod_r\lambda_{ir})I\) and \(ZS_i=\xi S_i Z\).
On the tensor product of the \(h\) modules take
\[
 L=\sum_{i=1}^h Z^{\otimes(i-1)}\otimes S_i\otimes I^{\otimes(h-i)} .
\]
Each summand \(V\) satisfies \(VU=\xi UV\) for \(U\) the sum of the
preceding summands. Potter's identity \cite{potter} gives
\((U+V)^b=U^b+V^b\) over the complex numbers. To see why it survives
specialization, order the factors in the expansion of \((U+V)^b\)
using \(VU=\xi UV\). The coefficient of \(U^{b-r}V^r\) is the
Gaussian binomial polynomial
\[
 \genfrac{[}{]}{0pt}{}{b}{r}_{\!X}
     =\prod_{j=1}^r\frac{1-X^{b-j+1}}{1-X^j}
\]
evaluated at \(X=\xi\), by the recurrence obtained on multiplying the
ordered expansion by \(U+V\). For \(0<r<b\) the denominators are
nonzero at a complex primitive \(b\)-th root, while the numerator
contains \(1-X^b\). Thus \(\Phi_b(X)\) divides the Gaussian binomial
in \(\mathbb Z[X]\). Every mixed coefficient vanishes over \(R_\ell[\xi]\)
and under any specialization. The \(b\)-th power of \(L\) is therefore
the sum of the \(b\)-th powers of its summands, which is \(pI\).
\end{proof}

Take \(b\) to be the least power of \(\ell\) at least \(\max(2,D)\).
Apply Lemma~\ref{lem:matrix-root} to every \(P_\beta\), padding the
monomial lists by zero terms to obtain matrices of one dimension.
Give each shift a block of that many independent rows. On the direct
sum of the blocks, \(b\) repetitions of one pencil stage compute a
predictor \(P\). At each frequency \(P\) is diagonal in the individual
rows. Write \(p_i\) for its multiplier on row \(i\). The exponent tuple
specifies the set \(J=J(\nu)\) of rows outside the window. It has size
at most \(s=\lfloor W/(2m^3)\rfloor\) and contains every row where
\(p_i\ne z\).

\subsection{Sending the error through a small spectral table}

At each frequency the prediction error has at most \(s\) known
positions. We can recover their values from \(s\) linear checks.
Returning them to the correct ordinary rows poses a second problem.
The set of destination rows changes with the frequency. Our lift
from the inverse-transformed rows uses a fixed row matrix, which acts
identically at every frequency. To make that lift work, we replace the
predictor by a polynomial whose error vanishes to second order at the
target multiplier.

\begin{lemma}[Full-rate Hermite repair]\label{lem:hermite-repair}
Work over a commutative ring in which \(\ell\) is invertible and there
are \(W\) elements with pairwise unit differences. Suppose that at each
frequency \(F\) has the same multiplier \(z\) on all rows, with
\(z^\ell=1\), and that \(P\) has known row multipliers \(p_i\) agreeing
with \(z\) off a known set \(J\) of size at most \(s\). Then \(F\) on
\(W\) independent inputs can be computed with at most \(\ell+2\)
applications of \(P\), one spectral table on \(2s\) rows and
pointwise linear maps.
\end{lemma}
\begin{proof}
If \(s=0\), \(P\) already gives \(F\). Otherwise, choose the \(s\times W\)
Vandermonde matrix \(H\) with powers \(0,\ldots,s-1\) on the given
nodes. Over a field these are the parity checks of a generalized
Reed--Solomon code \cite{reedsolomon}. For \(j=|J|\le s\), the first
\(j\) rows of the column submatrix \(H_J\) give a Vandermonde
determinant that is a unit. Hence \(H_J\) has a left inverse and
\(H_J^{\mathsf T}\) a right inverse.

On the batch \(X\), form \(PX\) and convert the \(2s\) check rows
\(HPX,HX\) to spectral coordinates. At one frequency the error vector
\(e\) satisfies
\[
 e_i=(p_i-z)\widehat X_i,\qquad
 He=H\widehat{PX}-zH\widehat X.
\]
The left inverse of \(H_J\) recovers \(e_J\) from \(He\).

To correct the output, first replace the predictor by the polynomial
\(g(P)\), where
\[
 g(y)=y-\ell^{-1}y(y^\ell-1).
\]
At every allowed \(z\), \(g(z)=z\) and \(g'(z)=0\). Dividing by the
monic polynomial \((y-z)^2\) gives
\(g(y)-z=(y-z)^2D_z(y)\) over the ring. Use the right inverse of
\(H_J^{\mathsf T}\) to choose an \(s\)-vector \(a\), linear in \(e_J\),
satisfying
\[
 (H^{\mathsf T}a)_i=D_z(p_i)e_i\qquad(i\in J).
\]
The table returns the pair \(a,za\). After inverse conversion and
pointwise multiplication by \(H^{\mathsf T}\), these are ordinary
rows \(Z,FZ\) with \(\widehat Z=H^{\mathsf T}a\)
(Figure~\ref{fig:repair-table}). Output \(g(P)X-PZ+FZ\).
Its error on row \(i\) at this frequency is
\[
 (p_i-z)\bigl(D_z(p_i)e_i-\widehat Z_i\bigr).
\]
The second factor vanishes on \(J\), and the first vanishes outside
\(J\). Thus the entire batch is correct. We reuse the initial \(PX\)
to form \(P^{\ell+1}X\) for \(g(P)X\), and apply \(P\) once more
to \(Z\). This totals \(\ell+2\) predictor applications.
\end{proof}

\begin{figure}[!htb]
\centering
\begin{tikzpicture}[
  font=\small,align=center,
  repair box/.style={draw=#1,fill=#1!4,rounded corners=1.5pt,
    inner xsep=6pt,inner ysep=4pt,text width=4.1cm,minimum height=1.26cm},
  repair path/.style={-{Stealth[length=1.7mm]},draw=ink,line width=.55pt},
  repair note/.style={font=\footnotesize,inner sep=2pt,align=center}
]
\node[anchor=south,font=\small\bfseries] at (0,.89) {Ordinary coordinates};
\node[anchor=south,font=\small\bfseries] at (7.7,.89) {At one frequency \(\nu\)};
\node[repair box=sea] (checks) at (0,0)
  {Form \(HX\) and \(HPX\),\\ each on \(s\) rows.};
\node[repair box=rust] (recover) at (7.7,0)
  {Use \(H\widehat{PX}-zH\widehat X\)\\ to recover \(e_J\).};
\node[repair box=rust] (coeffs) at (7.7,-1.95)
  {Choose \(a\) so \((H^{\mathsf T}a)_i=D_z(p_i)e_i\) on \(J\). Return \(a\) and \(za\).};
\node[repair box=sea] (lift) at (0,-1.95)
  {Lift by \(H^{\mathsf T}\) to obtain\\ \(Z\) and \(FZ\) on all \(W\) rows.};
\draw[repair path] (checks.east) -- node[repair note,above]{\(\mathcal S\)}
   node[repair note,below]{\(2s\) rows} (recover.west);
\draw[repair path] (recover.south) -- (coeffs.north);
\draw[repair path] (coeffs.west) -- node[repair note,above]{\(\mathcal S^{-1}\)}
  node[repair note,below]{\(2s\) rows} (lift.east);
\end{tikzpicture}
\caption{Repairing the predictor on \(W\) independent inputs. Only
\(2s\) rows cross into spectral coordinates and back. The table
prescribes the lifted correction on \(J(\nu)\). At every other row
\(p_i=z\), so subtracting \(PZ-FZ\) has no effect there.}
\label{fig:repair-table}
\end{figure}


\begin{proof}[Proof of Theorem~\ref{thm:repair}]
With \(W\) fixed, choose \(t\ge1\) so that \(\ell^t\ge W\) and
\(b\mid\ell^t\). Take a primitive \(\ell^t\)-th root \(\eta\) with
\(\eta^{\ell^{t-1}}=\zeta\). The cyclotomic degrees give the basis
\(1,\eta,\ldots,\eta^{\ell^{t-1}-1}\) for
\(R'=\mathbb Z[1/\ell,\eta]\) over \(R_\ell\). This ring contains the
matrix coefficients and supplies \(W\) Vandermonde nodes from the
powers of \(\eta\). For the quotient \(\rho\) of any two distinct
nodes, of order \(\ell^h\), the product of \(1-\rho^j\) over \(j\)
prime to \(\ell\) modulo \(\ell^h\) is \(\Phi_{\ell^h}(1)=\ell\).
Each factor is a unit, so the nodes have pairwise unit differences.

Each predictor application takes \(b\) pencil stages. Hence
Lemma~\ref{lem:hermite-repair} uses at most
\((\ell+2)b=O_\ell(\sqrt{m\log(2m)})\) stages and converts \(2s\)
rows in each direction. The bases, inverses and finite list of tables
can be constructed in \(R'\) from \(\ell^{-1}\) and the indicated
roots. These ring identities give the asserted program after base
change to every commutative \(R_\ell\)-algebra.
\end{proof}

The same construction applies to swaps of address fields. The following
form of the pencil is used for integer multiplication on tapes
\cite{oaT}.

\begin{corollary}[Simultaneous swaps in odd characteristic]
\label{cor:swap-repair}
Fix an odd prime \(p\) and a finite commutative ring \(K\) of
characteristic \(p\). Suppose \(Q_1,\ldots,Q_m\) interchange disjoint
pairs of address fields of equal respective lengths. For some fixed
batch width \(W\), their product on \(W\) independent arrays over \(K\)
can be applied with \(O(\sqrt{m\log(2m)})\) stages
\(A_0\otimes I+\sum_j A_j\otimes Q_j\) and one spectral table on
\(2s\le W/m^3\) rows. Forward or inverse conversion of one row uses
at most \(m\) individual swaps and address-dependent pointwise linear
maps.
\end{corollary}
\begin{proof}
Take \(\ell=2\) and \(\zeta=-1\). On a two-point orbit of one swap,
store the sum and difference in the two positions in a fixed order
determined by that pair of address fields. A fixed point keeps its
value. The two coordinates have eigenvalues \(+1,-1\), and the inverse
uses \(1/2\). Different swaps act on disjoint pairs, so these changes
of coordinates commute. A full conversion uses one swapped copy and
pointwise combinations for each pair. The resulting signs determine
the repair table.

Specializing the extension in Theorem~\ref{thm:repair} gives the finite
free \(K\)-algebra \(K'=K[Y]/(Y^{2^{t-1}}+1)\), of rank \(d=2^{t-1}\).
Pack \(d\) independent \(K\)-valued rows into one \(K'\)-valued row in
the monomial basis. Each prescribed scalar multiplication becomes a
fixed \(K\)-linear map on these \(d\) coordinates, while the swaps and
spectral conversions act componentwise. Expanding in this basis
multiplies both the batch width and the table width by \(d\), preserves
the number of pencil stages and returns the swap product on all the
original \(K\)-valued rows.
\end{proof}
